\chapter{Giới thiệu}
\label{chap:introduction}

\section{Bối cảnh nghiên cứu}
\label{sec:background}

Thực khuẩn thể (bacteriophage), thường được gọi tắt là phage, là các virus xâm nhiễm và nhân lên bên trong tế bào vi khuẩn. Với số lượng ước tính lên tới $10^{31}$ cá thể~\cite{suttle2007marine}, thực khuẩn thể là dạng sống phong phú nhất trên Trái Đất và đóng vai trò trung tâm trong việc điều hòa quần thể vi khuẩn, duy trì cân bằng sinh thái vi sinh vật, và thúc đẩy quá trình chuyển gen ngang (horizontal gene transfer) giữa các vi khuẩn~\cite{koskella2014bacteria}.

Thực khuẩn thể được phân loại thành hai nhóm chính dựa trên chiến lược nhân lên~\cite{wu2021deephage}. Thực khuẩn thể độc lực (virulent phage), còn gọi là phage phân giải nghiêm ngặt (strictly lytic phage), chỉ thực hiện chu trình phân giải: sau khi xâm nhập tế bào vi khuẩn, phage chiếm đoạt bộ máy tổng hợp của tế bào chủ để sản xuất các virion con, rồi gây phá vỡ tế bào để giải phóng thế hệ phage mới~\cite{mcnair2012phacts}. Ngược lại, thực khuẩn thể ôn hòa (temperate phage) có thể thực hiện cả chu trình phân giải lẫn chu trình lysogen: trong chu trình lysogen, vật liệu di truyền của phage được tích hợp vào nhiễm sắc thể vi khuẩn chủ dưới dạng prophage, tồn tại ở trạng thái không hoạt động và nhân lên cùng hệ gen vi khuẩn qua nhiều thế hệ~\cite{howard2017lysogeny}. Prophage có thể được kích hoạt bởi các yếu tố môi trường như tia UV hoặc tác nhân gây tổn thương DNA, chuyển sang chu trình phân giải~\cite{feiner2015new}.

Việc phân loại chính xác lối sống của thực khuẩn thể có ý nghĩa quan trọng trong cả nghiên cứu vi sinh vật học và ứng dụng y sinh học. Trong nghiên cứu hệ vi sinh vật, xác định lối sống phage giúp làm sáng tỏ các tương tác phức tạp giữa phage và vi khuẩn chủ, từ đó hiểu rõ hơn vai trò sinh thái của phage trong các cộng đồng vi sinh vật~\cite{shang2023phatyp}. Trong bối cảnh kháng kháng sinh ngày càng gia tăng, liệu pháp phage (phage therapy) đang được xem xét như một giải pháp thay thế đầy hứa hẹn cho kháng sinh truyền thống~\cite{gorski2016phage,azimi2019phage}. Tuy nhiên, chỉ có thực khuẩn thể độc lực mới phù hợp cho liệu pháp phage, vì thực khuẩn thể ôn hòa có thể tích hợp vào hệ gen vi khuẩn bệnh nguyên và tiềm ẩn nguy cơ truyền các gen độc lực hoặc gen kháng kháng sinh thông qua quá trình lysogenic conversion~\cite{cobian2016viruses}.

Các phương pháp truyền thống để xác định lối sống phage dựa vào nuôi cấy trong phòng thí nghiệm và quan sát kiểu hình, nhưng tốn thời gian, đòi hỏi nhiều công sức, và không thể áp dụng cho phần lớn các phage chưa được nuôi cấy được trong các mẫu môi trường~\cite{shang2023phatyp}. Những hạn chế này đã thúc đẩy sự phát triển của các phương pháp tính toán để dự đoán lối sống phage trực tiếp từ dữ liệu trình tự gen. Với sự phát triển của công nghệ giải trình tự thế hệ mới, metagenomics đã trở thành công cụ quan trọng để nghiên cứu đa dạng sinh học của phage mà không cần nuôi cấy~\cite{mokili2012metagenomics,hayes2017metagenomic}. Trong phân tích metagenomic, hệ gen phage thường được phục hồi dưới dạng các contig—các chuỗi liên tục được lắp ráp từ các đoạn đọc ngắn chồng lấp~\cite{setubal2021metagenome,ghurye2016metagenomic}. Các contig này thường có độ dài biến thiên lớn và bị phân mảnh do độ phủ giải trình tự không đồng đều, các hiện vật kỹ thuật, và các vùng gen lặp lại, khiến việc phân loại chính xác các chuỗi không hoàn chỉnh trở thành một thách thức tính toán quan trọng.

\section{Vấn đề nghiên cứu}
\label{sec:problem}

Các phương pháp tính toán để dự đoán lối sống thực khuẩn thể phải đối mặt với bốn thách thức chính. Thứ nhất, cơ sở dữ liệu hệ gen phage hiện có vẫn còn rất hạn chế so với sự đa dạng khổng lồ của phage trong tự nhiên, phần lớn chưa được phân lập và nghiên cứu đặc tính~\cite{hayes2017metagenomic}. Thứ hai, không giống vi khuẩn có gen 16S rRNA làm gen đánh dấu phổ quát, thực khuẩn thể không có gen đánh dấu tương tự~\cite{mokili2012metagenomics}, khiến các phương pháp phân loại dựa trên gen đánh dấu không khả thi. Thứ ba, dữ liệu metagenomic thường chứa nhiều đoạn DNA ngắn bao gồm chỉ một phần gen hoặc các chuỗi gen không hoàn chỉnh, làm phức tạp các chiến lược phân loại dựa trên gen chức năng~\cite{jiang2017ensvmb}. Thứ tư, thực khuẩn thể có thể chia sẻ các chuỗi gen với vi khuẩn chủ thông qua chuyển gen ngang, đặc biệt trong trường hợp prophage được tích hợp vào hệ gen vi khuẩn~\cite{rozov2017recycler}, dẫn đến phân loại sai trong quá trình gán phân loại học.

Dựa trên các thách thức trên, bài toán dự đoán lối sống thực khuẩn thể được phát biểu như sau: cho một tập hợp các contig $\mathcal{S} = \{S_1, S_2, \ldots, S_n\}$ thu được từ giải trình tự và lắp ráp metagenomic, trong đó mỗi contig $S_i$ có độ dài $L_i$ base pairs, nhiệm vụ là học hàm phân loại $f: \mathcal{S} \rightarrow \{0, 1\}$ trong đó $0$ đại diện cho lớp temperate và $1$ đại diện cho lớp virulent, sao cho $f$ hoạt động hiệu quả trên các contig có độ dài biến thiên lớn mà không phụ thuộc vào các công cụ dự đoán gen bên ngoài.

Các phương pháp hiện có đã đạt được tiến bộ đáng kể nhưng vẫn còn những hạn chế quan trọng. Các phương pháp học máy truyền thống như PHACTS~\cite{mcnair2012phacts}, PhageAI~\cite{tynecki2020phageai}, và BACPHLIP~\cite{hockenberry2021bacphlip} được thiết kế cho hệ gen hoàn chỉnh và hiệu suất giảm mạnh trên contig phân mảnh. DeePhage~\cite{wu2021deephage} sử dụng CNN trên mã hóa one-hot nhưng thiếu tri thức tiền huấn luyện, hạn chế khả năng tổng quát hóa trên các contig ngắn với thông tin ngữ cảnh hạn chế. PhaTYP~\cite{shang2023phatyp} cải thiện hiệu suất nhờ kiến trúc BERT nhưng phụ thuộc vào Prodigal~\cite{hyatt2010prodigal} để dự đoán protein—công cụ này có hiệu suất giảm đáng kể trên các đoạn ngắn hơn 200 bp~\cite{trimble2012short}, khiến PhaTYP kém hiệu quả trên contig ngắn thiếu vùng mã hóa protein hoàn chỉnh. DeepPL~\cite{zhang2024deeppl} và ProkBERT~\cite{ligeti2024prokbert} sử dụng tokenization k-mer chồng lấp, gây ra rò rỉ thông tin trong quá trình masked language modeling và kém hiệu quả tính toán do độ dài chuỗi token gần bằng độ dài chuỗi gốc~\cite{zhou2023dnabert}.

\section{Mục tiêu và đóng góp}
\label{sec:objectives}

Mục tiêu chính của luận văn là phát triển PhaBERT-CNN, một phương pháp học sâu lai kết hợp mô hình nền tảng genome DNABERT-2~\cite{zhou2023dnabert} với mạng nơ-ron tích chập đa tỷ lệ và cơ chế attention pooling để phân loại lối sống thực khuẩn thể từ các contig metagenomic. Phương pháp được thiết kế để hoạt động trực tiếp trên chuỗi DNA thô thông qua BPE tokenization, không phụ thuộc vào công cụ dự đoán gene, và được đánh giá toàn diện trên bốn nhóm độ dài contig (100--1800 bp) thông qua stratified 5-fold cross-validation.

Luận văn đóng góp vào lĩnh vực sinh tin học thông qua bốn đóng góp chính. Thứ nhất, đề xuất kiến trúc PhaBERT-CNN mới kết hợp DNABERT-2 với mô đun Modified TextCNN gồm nhánh CNN đa tỷ lệ (kernel sizes 3, 5, 7) và nhánh attention pooling xử lý song song, tạo ra biểu diễn tổng hợp 512 chiều kết hợp cả đặc trưng cục bộ đa tỷ lệ và ngữ cảnh toàn cục. Thứ hai, đề xuất chiến lược huấn luyện hai giai đoạn với discriminative learning rates: giai đoạn warm-up đóng băng backbone để ổn định các lớp task-specific, tiếp theo là full fine-tuning với learning rate thấp hơn cho DNABERT-2 ($1 \times 10^{-5}$) và cao hơn cho các lớp task-specific ($1 \times 10^{-4}$). Thứ ba, thực hiện đánh giá có hệ thống so sánh PhaBERT-CNN với bốn phương pháp tiên tiến (DNABERT-2, PhaTYP, DeePhage, ProkBERT) trên bốn nhóm độ dài contig, chứng minh hiệu suất vượt trội trên contig ngắn và trung bình (accuracy 81.59\%--90.01\% trên Nhóm A--C) và kết quả cạnh tranh trên contig dài (90.69\% ở Nhóm D). Thứ tư, cung cấp phân tích chi tiết về ưu điểm, hạn chế, và hướng phát triển của phương pháp, đặt nền móng cho các nghiên cứu tương lai về ứng dụng genomic foundation models vào các bài toán sinh tin học downstream.

\section{Phạm vi nghiên cứu}
\label{sec:scope}

Luận văn tập trung vào bài toán phân loại nhị phân lối sống thực khuẩn thể (virulent/temperate) từ các contig DNA trong phạm vi độ dài 100--1800 bp, được chia thành bốn nhóm: Nhóm A (100--400 bp), Nhóm B (400--800 bp), Nhóm C (800--1200 bp), và Nhóm D (1200--1800 bp). Tập dữ liệu được xây dựng từ 2,241 hệ gen phage hoàn chỉnh (707 temperate, 1,534 virulent) tổng hợp từ DeePhage~\cite{wu2021deephage} và DeepPL~\cite{zhang2024deeppl}; các contig được tạo bằng phương pháp sliding window để mô phỏng kịch bản metagenomic thực tế. Mô hình nền tảng sử dụng là DNABERT-2~\cite{zhou2023dnabert}; các mô hình nền tảng khác không được đánh giá trong phạm vi luận văn này. Hiệu suất được đo bằng ba chỉ số sensitivity, specificity, và accuracy nhất quán với các nghiên cứu trước, sử dụng stratified 5-fold cross-validation. Các contig đầu vào được giả định đã được xác định là có nguồn gốc từ phage; luận văn không xem xét bước nhận dạng phage từ dữ liệu metagenomic thô.

Luận văn tập trung vào phạm vi nghiên cứu cụ thể sau đây, với các giới hạn và giả định rõ ràng:

\subsection{Phạm vi về bài toán}

\textbf{Phân loại nhị phân lối sống:} Luận văn tập trung vào bài toán phân loại nhị phân (binary classification) để phân biệt thực khuẩn thể độc lực (virulent/lytic) và thực khuẩn thể ôn hòa (temperate/lysogenic). Các lối sống phức tạp hơn hoặc các trạng thái trung gian không được xem xét trong phạm vi nghiên cứu này.

\textbf{Dữ liệu đầu vào:} Phương pháp được thiết kế để hoạt động trên các contig DNA được lắp ráp từ dữ liệu metagenomic, không phải trên các đoạn đọc thô (raw reads). Các contig được giả định là đã được xác định là có nguồn gốc từ phage thông qua các công cụ nhận dạng phage như VirSorter, VIBRANT, hoặc các phương pháp tương tự.

\subsection{Phạm vi về dữ liệu}

\textbf{Phạm vi độ dài contig:} Nghiên cứu tập trung vào bốn nhóm độ dài contig:
\begin{itemize}
    \item Nhóm A: 100-400 bp (contig rất ngắn)
    \item Nhóm B: 400-800 bp (contig ngắn)
    \item Nhóm C: 800-1200 bp (contig trung bình)
    \item Nhóm D: 1200-1800 bp (contig dài)
\end{itemize}

Các contig ngắn hơn 100 bp hoặc dài hơn 1800 bp không được đánh giá chi tiết trong nghiên cứu này, mặc dù phương pháp về mặt kỹ thuật có thể xử lý các chuỗi có độ dài tùy ý nhờ ALiBi positional encoding của DNABERT-2.

\textbf{Nguồn dữ liệu:} Tập dữ liệu được xây dựng từ hai nghiên cứu gần đây: DeePhage \cite{wu2021deephage} và DeepPL \cite{zhang2024deeppl}, bao gồm 2,241 hệ gen phage hoàn chỉnh (707 temperate, 1,534 virulent). Các contig được tạo ra bằng phương pháp sliding window để mô phỏng kịch bản metagenomic thực tế.

\textbf{Giới hạn về đa dạng sinh học:} Mặc dù tập dữ liệu bao gồm nhiều họ phage khác nhau, nó vẫn chỉ đại diện cho một phần nhỏ của sự đa dạng phage trong tự nhiên. Hiệu suất của mô hình trên các phage thuộc các họ hiếm gặp hoặc chưa được nghiên cứu có thể khác với kết quả báo cáo.

\subsection{Phạm vi về phương pháp}

\textbf{Mô hình nền tảng:} Nghiên cứu sử dụng DNABERT-2 \cite{zhou2023dnabert} làm mô hình nền tảng genome. Các mô hình nền tảng khác như Nucleotide Transformer \cite{dalla2025nucleotide}, HyenaDNA, hoặc các biến thể BERT khác không được đánh giá trong phạm vi luận văn này.

\textbf{Kiến trúc downstream:} Luận văn tập trung vào kiến trúc lai kết hợp CNN đa tỷ lệ và attention pooling. Các kiến trúc downstream khác như RNN, LSTM, hoặc pure transformer decoder không được khảo sát chi tiết.

\textbf{Phương pháp so sánh:} Luận văn so sánh với bốn phương pháp tiên tiến: DNABERT-2, PhaTYP, DeePhage, và ProkBERT. Các phương pháp học máy truyền thống như PHACTS, PhageAI, và BACPHLIP không được đánh giá lại vì đã được chứng minh có hiệu suất thấp hơn đáng kể trên dữ liệu contig trong các nghiên cứu trước.

\subsection{Phạm vi về đánh giá}

\textbf{Metrics đánh giá:} Luận văn sử dụng ba metrics chính: sensitivity (recall), specificity, và accuracy. Các metrics khác như precision, F1-score, AUC-ROC, hoặc Matthews correlation coefficient không được báo cáo chi tiết để duy trì tính nhất quán với các nghiên cứu trước đó.

\textbf{Phương pháp validation:} Sử dụng stratified 5-fold cross-validation để đánh giá hiệu suất. Không thực hiện đánh giá trên các tập dữ liệu độc lập bên ngoài (external validation) do hạn chế về tài nguyên và tính sẵn có của dữ liệu có nhãn chất lượng cao.

\textbf{Phân tích ablation:} Luận văn thực hiện phân tích so sánh giữa PhaBERT-CNN và DNABERT-2 baseline để đánh giá tác động của các thành phần CNN và attention. Các nghiên cứu ablation chi tiết hơn (ví dụ: loại bỏ từng nhánh CNN, thay đổi kernel sizes, so sánh các cơ chế attention khác nhau) không được thực hiện trong phạm vi luận văn này.

\subsection{Phạm vi về triển khai}

\textbf{Framework và thư viện:} Mô hình được triển khai sử dụng PyTorch và thư viện Transformers của Hugging Face. Các framework khác như TensorFlow/Keras không được xem xét.

\textbf{Tài nguyên tính toán:} Các thí nghiệm được thực hiện trên GPU NVIDIA RTX A4000 16GB. Phân tích về khả năng mở rộng (scalability) trên các cấu hình phần cứng khác hoặc tối ưu hóa cho triển khai trên CPU không được đề cập chi tiết.

\textbf{Công cụ phần mềm:} Luận văn không phát triển một công cụ phần mềm độc lập (standalone software tool) hoặc web server. Mã nguồn được cung cấp dưới dạng script Python để tái tạo kết quả nghiên cứu.

\subsection{Giới hạn và giả định}

\textbf{Giả định về chất lượng dữ liệu:} Các contig đầu vào được giả định là:
\begin{itemize}
    \item Đã được xác định chính xác là có nguồn gốc từ phage
    \item Không chứa lỗi giải trình tự đáng kể
    \item Không bị ô nhiễm bởi DNA vi khuẩn chủ
\end{itemize}

\textbf{Giới hạn về tổng quát hóa:} Hiệu suất của mô hình được đánh giá trên dữ liệu mô phỏng (simulated data) được tạo từ hệ gen hoàn chỉnh. Hiệu suất trên dữ liệu metagenomic thực tế từ các môi trường đa dạng có thể khác biệt và cần được xác thực thêm.

\textbf{Không xem xét các yếu tố sinh học phức tạp:} Luận văn không xem xét các trường hợp đặc biệt như:
\begin{itemize}
    \item Phage có khả năng chuyển đổi giữa các lối sống (lifestyle switching)
    \item Prophage bị đột biến mất chức năng (defective prophage)
    \item Phage với các cơ chế lysogen hóa không điển hình
    \item Tác động của môi trường chủ lên biểu hiện lối sống
\end{itemize}

Các giới hạn và giả định này được thiết lập để đảm bảo nghiên cứu tập trung và khả thi trong phạm vi luận văn thạc sĩ, đồng thời cung cấp nền tảng vững chắc cho các nghiên cứu mở rộng trong tương lai.

\section{Bố cục luận văn}
\label{sec:organization}

Luận văn được tổ chức thành năm chương. \textbf{Chương 1} giới thiệu bối cảnh nghiên cứu về thực khuẩn thể, các thách thức trong phân loại lối sống phage từ dữ liệu metagenomic, mục tiêu và đóng góp chính, phạm vi nghiên cứu, và bố cục luận văn. \textbf{Chương 2} cung cấp nền tảng lý thuyết về sinh học thực khuẩn thể, metagenomics, các kiến trúc học sâu (CNN, Transformer, BERT), mô hình DNABERT-2, và tổng quan các công trình liên quan về phân loại lối sống phage. \textbf{Chương 3} trình bày chi tiết phương pháp PhaBERT-CNN, bao gồm quy trình xây dựng dữ liệu (sliding window, reverse complement augmentation, Random Undersampling), kiến trúc mô hình (DNABERT-2 backbone, Modified TextCNN với nhánh CNN đa tỷ lệ và nhánh attention pooling, đầu phân loại), và chiến lược huấn luyện hai giai đoạn với discriminative learning rates. \textbf{Chương 4} trình bày thiết lập thực nghiệm, kết quả so sánh PhaBERT-CNN với DNABERT-2 baseline và các phương pháp tiên tiến (PhaTYP, DeePhage, ProkBERT) trên bốn nhóm độ dài contig. \textbf{Chương 5} tổng kết đóng góp của luận văn, phân tích ưu điểm và hạn chế của phương pháp dựa trên kết quả thực nghiệm, và đề xuất các hướng nghiên cứu tương lai.
Luận văn được tổ chức thành năm chương với cấu trúc logic như sau:

\textbf{Chương 1: Giới thiệu}

Chương này giới thiệu bối cảnh nghiên cứu về thực khuẩn thể và tầm quan trọng của việc phân loại lối sống phage trong nghiên cứu vi sinh vật học và ứng dụng y sinh học. Chương trình bày các thách thức chính trong phân loại lối sống phage từ dữ liệu metagenomic, bao gồm hạn chế về dữ liệu, thiếu gen đánh dấu phổ quát, độ dài đoạn DNA ngắn, và vấn đề tương đồng trình tự. Chương cũng phát biểu rõ ràng bài toán nghiên cứu, mục tiêu, đóng góp chính, và phạm vi của luận văn.

\textbf{Chương 2: Cơ sở lý thuyết}

Chương này cung cấp nền tảng lý thuyết toàn diện cần thiết để hiểu phương pháp đề xuất. Chương bắt đầu với kiến thức sinh học về thực khuẩn thể, bao gồm phân loại, chu trình sống, và vai trò sinh thái. Tiếp theo, chương giới thiệu các khái niệm về metagenomics và giải trình tự thế hệ mới, giải thích cách các contig được tạo ra từ dữ liệu metagenomic. Phần tiếp theo trình bày các nền tảng học máy và học sâu, bao gồm mạng nơ-ron tích chập (CNN), kiến trúc transformer, và cơ chế attention. Chương dành phần quan trọng để giải thích chi tiết về BERT và các biến thể của nó cho dữ liệu genome (DNABERT, DNABERT-2), bao gồm masked language modeling, tokenization strategies, và transfer learning. Cuối cùng, chương tổng hợp các công trình liên quan về phân loại lối sống phage, phân tích ưu nhược điểm của các phương pháp học máy truyền thống (PHACTS, PhageAI, BACPHLIP) và các phương pháp học sâu hiện đại (DeePhage, PhaTYP, DeepPL, ProkBERT).

\textbf{Chương 3: Phương pháp đề xuất}

Chương này trình bày chi tiết phương pháp PhaBERT-CNN được đề xuất. Chương bắt đầu với tổng quan về kiến trúc tổng thể và động lực thiết kế. Tiếp theo, chương mô tả quy trình xây dựng tập dữ liệu, bao gồm nguồn dữ liệu, phương pháp tạo contig bằng sliding window, data augmentation thông qua reverse complement, và chiến lược xử lý mất cân bằng lớp. Chương sau đó giải thích chi tiết về mô hình nền tảng DNABERT-2, bao gồm chiến lược tokenization BPE, các cải tiến kiến trúc (ALiBi, Flash Attention), và corpus tiền huấn luyện đa loài. Phần cốt lõi của chương mô tả kiến trúc PhaBERT-CNN với hai nhánh trích xuất đặc trưng song song: (1) multi-scale CNN feature extractor với ba kernel sizes khác nhau, và (2) attention-based global feature extractor. Chương kết thúc bằng việc trình bày chiến lược huấn luyện hai giai đoạn với learning rate phân biệt, bao gồm các chi tiết về optimizer, learning rate schedule, và early stopping.

\textbf{Chương 4: Thực nghiệm và đánh giá}

Chương này trình bày thiết lập thí nghiệm toàn diện và kết quả đánh giá. Chương bắt đầu với mô tả chi tiết về môi trường thực nghiệm, bao gồm cấu hình phần cứng, phần mềm, và các hyperparameters. Tiếp theo, chương giải thích các metrics đánh giá (sensitivity, specificity, accuracy) và phương pháp validation (stratified 5-fold cross-validation). Chương trình bày kết quả thực nghiệm theo hai phần chính: (1) so sánh PhaBERT-CNN với DNABERT-2 baseline để đánh giá tác động của các thành phần CNN và attention, và (2) so sánh với các phương pháp tiên tiến (PhaTYP, DeePhage, ProkBERT) trên bốn nhóm độ dài contig. Kết quả được trình bày thông qua các bảng số liệu chi tiết và biểu đồ trực quan. Chương cũng bao gồm phân tích chi tiết về hiệu suất trên từng nhóm độ dài, phân tích lỗi, và thảo luận về các trường hợp mô hình hoạt động tốt và các trường hợp thất bại. Cuối cùng, chương thảo luận về trade-off giữa độ chính xác và thời gian tính toán, cũng như các hạn chế của phương pháp đề xuất.

\textbf{Chương 5: Kết luận}

Chương cuối cùng tổng kết các kết quả chính của luận văn và đóng góp của nghiên cứu. Chương nhắc lại bài toán nghiên cứu và tóm tắt cách PhaBERT-CNN giải quyết các thách thức đã xác định. Chương tổng hợp các phát hiện quan trọng từ thực nghiệm, bao gồm hiệu suất vượt trội trên các contig ngắn và trung bình, sự cân bằng giữa sensitivity và specificity, và so sánh với các phương pháp hiện có. Chương thảo luận về các hạn chế của nghiên cứu, bao gồm phụ thuộc vào chất lượng dữ liệu đầu vào, nhu cầu về tài nguyên tính toán, và các giả định về dữ liệu. Cuối cùng, chương đề xuất các hướng nghiên cứu tương lai, bao gồm: (1) đánh giá trên dữ liệu metagenomic thực tế từ các môi trường đa dạng, (2) khảo sát các mô hình nền tảng genome khác, (3) phát triển các chiến lược xử lý mất cân bằng lớp hiệu quả hơn, (4) tối ưu hóa cho triển khai thực tế, và (5) mở rộng sang các nhiệm vụ phân loại phage khác như phân loại họ phage hoặc dự đoán vi khuẩn chủ.

\vspace{1em}

Cấu trúc này được thiết kế để dẫn dắt người đọc từ động lực và nền tảng lý thuyết, qua phương pháp đề xuất chi tiết, đến đánh giá thực nghiệm toàn diện, và cuối cùng là kết luận và hướng phát triển. Mỗi chương xây dựng dựa trên các chương trước, tạo thành một câu chuyện nghiên cứu mạch lạc và logic.
